% =====================================================================
%  BUỔI 16 — KẾ HOẠCH GIẢNG DẠY (từ giao-vien/01-ke-hoach-chua-de.md)
%  Sinh tự động bằng scripts/md2tex-chung.py — không sửa tay.
% =====================================================================
\documentclass[11pt, a4paper]{article}
% =====================================================================
%  CAE SHANGHAI — GÓI LỆNH DÙNG CHUNG CHO TÀI LIỆU LUYỆN THI CSCA TOÁN
%  Biên dịch bằng XeLaTeX. Mọi file .tex trong buổi học \input file này.
% =====================================================================

% ---------- NGÔN NGỮ & PHÔNG CHỮ ----------
\usepackage{fontspec}
\usepackage{polyglossia}
\setmainlanguage{vietnamese}
\setotherlanguage{english}

% Chữ Latin (Anh + Việt): TeX Gyre Termes — bản clone Times New Roman,
% đi kèm TeX Live, hỗ trợ đầy đủ dấu tiếng Việt.
\setmainfont{TeX Gyre Termes}
\setsansfont{TeX Gyre Heros}

% Chữ Hán: ưu tiên Source Han Serif SC (思源宋体) — chuẩn học thuật Trung Quốc.
% Nếu máy chưa cài, tự động dùng Songti SC (宋体, có sẵn trên macOS).
% Cả hai đều thuộc họ 宋体 nên diện mạo tài liệu không đổi.
\IfFontExistsTF{Source Han Serif SC}{
  \newcommand{\cjkfamilyname}{Source Han Serif SC}
}{
  \IfFontExistsTF{Songti SC}{
    \newcommand{\cjkfamilyname}{Songti SC}
  }{
    \newcommand{\cjkfamilyname}{Noto Serif CJK SC}
  }
}

% xeCJK tự chuyển font khi gặp ký tự Hán, không cần bọc \cjk{} thủ công.
% Đây là điều kiện bắt buộc: font Latin (TeX Gyre Termes) không có glyph chữ Hán,
% nếu không khai báo xeCJK thì mọi chữ Hán viết trần sẽ bị nuốt khỏi bản in.
\usepackage{xeCJK}
\setCJKfamilyfont{caehan}{\cjkfamilyname}[Scale=0.95]
\setCJKmainfont{\cjkfamilyname}[Scale=0.95]
\setCJKsansfont{\cjkfamilyname}[Scale=0.95]
\setCJKmonofont{\cjkfamilyname}[Scale=0.95]

% \cjk{} ép dùng font Hán trong ngữ cảnh cần thiết (ví dụ trong tiêu đề in đậm
% hoặc trong bảng có định dạng riêng). Dùng \CJKfamily chứ không phải tên font,
% vì tên font chỉ là văn bản — sẽ in ra chữ thay vì đổi font.
\newcommand{\cjk}[1]{{\CJKfamily{caehan}#1}}

% ---------- TOÁN HỌC ----------
\usepackage{amsmath, amssymb, amsthm, mathtools}
\usepackage{unicode-math}
\setmathfont{TeX Gyre Termes Math}

% ---------- BỐ CỤC TRANG ----------
% Lề phải chừa vùng an toàn của letterhead:
%   banner trên chiếm 0–19mm  → top tối thiểu 2,2cm
%   banner dưới chiếm 281–297mm → bottom tối thiểu 2,2cm
% Lề trên 2,6cm (26mm) để tiêu đề đầu trang không đè banner (banner hết ở 19mm).
\usepackage[a4paper, top=2.6cm, bottom=2.4cm, left=2.5cm, right=2.0cm]{geometry}
\usepackage{fancyhdr}
\usepackage{titlesec}
\usepackage{enumitem}
\usepackage{graphicx}
\usepackage{booktabs, array, longtable, tabularx}
\usepackage{xcolor}
\usepackage[most]{tcolorbox}
\usepackage[colorlinks=true, linkcolor=cae_blue, citecolor=cae_blue, urlcolor=cae_blue]{hyperref}

% ---------- BẢNG MÀU THƯƠNG HIỆU ----------
\definecolor{cae_blue}{RGB}{0, 51, 102}
\definecolor{cae_gold}{RGB}{196, 149, 60}
\definecolor{box_bg}{RGB}{245, 247, 250}
\definecolor{box_line}{RGB}{0, 51, 102}

% ---------- TIÊU ĐỀ MỤC ----------
% Đệm đầu trang: banner letterhead chiếm 0–19mm, mà lề trên là 22mm.
% \section* đầu trang (sau \clearpage) sát lề nên phải chừa thêm, nếu không
% sẽ đè lên banner.
\titleformat{\section}
  {\normalfont\large\bfseries\sffamily\color{cae_blue}}{\thesection.}{0.6em}{}
\titleformat{\subsection}
  {\normalfont\normalsize\bfseries\sffamily\color{cae_blue!85!black}}{\thesubsection}{0.6em}{}
\titlespacing*{\section}{0pt}{14pt}{6pt}
\titlespacing*{\subsection}{0pt}{10pt}{4pt}

% Tiêu đề không đánh số ở đầu trang: \section* sau \clearpage nằm sát lề trên
% nên đè lên banner letterhead. Dùng \vspace* ngay TRƯỚC \section* không đủ,
% vì titlesec chèn khoảng cách riêng. Cách chắc chắn là đệm bằng \null.
\newcommand{\sectiondautrang}[1]{\section*{#1}}

% Dạng bài dùng \subsection*{} nên cần định dạng riêng, đậm và có gạch chân màu
\titleformat{\subsubsection}
  {\normalfont\normalsize\bfseries\sffamily\color{cae_blue!85!black}}{\thesubsubsection}{0.6em}{}

% ---------- LETTERHEAD (ảnh nền toàn trang) ----------
% Ảnh letterhead của CAE Shanghai: banner xanh trên/dưới kèm logo chìm giữa.
% Tỉ lệ ảnh 1414x2000 = 0.7070, khớp A4 (0.7071) nên chèn thẳng không méo.
% Đặt ở lớp nền dưới cùng để chữ không bị ảnh che.
% Dùng \AddToShipoutPictureBG (không có dấu *) để áp dụng cho MỌI trang.
\usepackage{eso-pic}
\AddToShipoutPictureBG{%
  \AtPageLowerLeft{%
    \includegraphics[width=\paperwidth,height=\paperheight]{images/letterhead.png}%
  }%
}

% ---------- ĐẦU / CHÂN TRANG ----------
% Letterhead đã có banner trên và dưới, nên bỏ dòng chạy để tránh chồng chữ.
% Số trang đặt trong vùng nội dung, sát trên banner dưới, lề phải.
\pagestyle{fancy}
\setlength{\headheight}{0pt}
\fancyhf{}
\renewcommand{\headrulewidth}{0pt}
\renewcommand{\footrulewidth}{0pt}
\fancyfoot[R]{\raisebox{0.9cm}{\small\color{cae_blue}\thepage}}

% ---------- MÔI TRƯỜNG NỘI DUNG ----------
% Trong tcolorbox, LaTeX đặt parindent = parskip = 0, nên các đoạn văn trong
% khung dính liền nhau. Đặt lại khoảng cách đoạn cho mọi khung.
\newcommand{\khoangcachdoan}{\setlength{\parskip}{0.5em}\setlength{\parindent}{0pt}}
\tcbset{before upper={\khoangcachdoan}}

\newtcolorbox{dinhnghia}[1][]{
  colback=box_bg, colframe=box_line, boxrule=0.6pt, arc=1pt,
  fonttitle=\bfseries, title={Định nghĩa #1}, breakable}
\newtcolorbox{dinhly}[1][]{
  colback=box_bg, colframe=cae_gold, boxrule=0.6pt, arc=1pt,
  fonttitle=\bfseries, title={Tính chất #1}, breakable}
\newtcolorbox{vidu}[1][]{
  colback=white, colframe=black!40, boxrule=0.4pt, arc=1pt,
  fonttitle=\bfseries, title={Ví dụ #1}, breakable}
\newtcolorbox{baitap}[1][]{
  colback=white, colframe=cae_blue!50, boxrule=0.4pt, arc=1pt,
  fonttitle=\bfseries, title={Câu #1}, breakable}

% Ô đáp án: viền xanh đậm hơn, dùng trong file đáp án
\newtcolorbox{dapan}[1][]{
  colback=box_bg, colframe=cae_blue, boxrule=0.6pt, arc=1pt,
  fonttitle=\bfseries, title={Câu #1}, breakable}

% ---------- TÊN CHỨNG MINH ----------
\renewcommand{\proofname}{\normalfont\bfseries Chứng minh}

% ---------- MÔN HỌC ----------
% Đổi được ở file .tex trước khi % =====================================================================
%  CAE SHANGHAI — GÓI LỆNH DÙNG CHUNG CHO TÀI LIỆU LUYỆN THI CSCA TOÁN
%  Biên dịch bằng XeLaTeX. Mọi file .tex trong buổi học \input file này.
% =====================================================================

% ---------- NGÔN NGỮ & PHÔNG CHỮ ----------
\usepackage{fontspec}
\usepackage{polyglossia}
\setmainlanguage{vietnamese}
\setotherlanguage{english}

% Chữ Latin (Anh + Việt): TeX Gyre Termes — bản clone Times New Roman,
% đi kèm TeX Live, hỗ trợ đầy đủ dấu tiếng Việt.
\setmainfont{TeX Gyre Termes}
\setsansfont{TeX Gyre Heros}

% Chữ Hán: ưu tiên Source Han Serif SC (思源宋体) — chuẩn học thuật Trung Quốc.
% Nếu máy chưa cài, tự động dùng Songti SC (宋体, có sẵn trên macOS).
% Cả hai đều thuộc họ 宋体 nên diện mạo tài liệu không đổi.
\IfFontExistsTF{Source Han Serif SC}{
  \newcommand{\cjkfamilyname}{Source Han Serif SC}
}{
  \IfFontExistsTF{Songti SC}{
    \newcommand{\cjkfamilyname}{Songti SC}
  }{
    \newcommand{\cjkfamilyname}{Noto Serif CJK SC}
  }
}

% xeCJK tự chuyển font khi gặp ký tự Hán, không cần bọc \cjk{} thủ công.
% Đây là điều kiện bắt buộc: font Latin (TeX Gyre Termes) không có glyph chữ Hán,
% nếu không khai báo xeCJK thì mọi chữ Hán viết trần sẽ bị nuốt khỏi bản in.
\usepackage{xeCJK}
\setCJKfamilyfont{caehan}{\cjkfamilyname}[Scale=0.95]
\setCJKmainfont{\cjkfamilyname}[Scale=0.95]
\setCJKsansfont{\cjkfamilyname}[Scale=0.95]
\setCJKmonofont{\cjkfamilyname}[Scale=0.95]

% \cjk{} ép dùng font Hán trong ngữ cảnh cần thiết (ví dụ trong tiêu đề in đậm
% hoặc trong bảng có định dạng riêng). Dùng \CJKfamily chứ không phải tên font,
% vì tên font chỉ là văn bản — sẽ in ra chữ thay vì đổi font.
\newcommand{\cjk}[1]{{\CJKfamily{caehan}#1}}

% ---------- TOÁN HỌC ----------
\usepackage{amsmath, amssymb, amsthm, mathtools}
\usepackage{unicode-math}
\setmathfont{TeX Gyre Termes Math}

% ---------- BỐ CỤC TRANG ----------
% Lề phải chừa vùng an toàn của letterhead:
%   banner trên chiếm 0–19mm  → top tối thiểu 2,2cm
%   banner dưới chiếm 281–297mm → bottom tối thiểu 2,2cm
% Lề trên 2,6cm (26mm) để tiêu đề đầu trang không đè banner (banner hết ở 19mm).
\usepackage[a4paper, top=2.6cm, bottom=2.4cm, left=2.5cm, right=2.0cm]{geometry}
\usepackage{fancyhdr}
\usepackage{titlesec}
\usepackage{enumitem}
\usepackage{graphicx}
\usepackage{booktabs, array, longtable, tabularx}
\usepackage{xcolor}
\usepackage[most]{tcolorbox}
\usepackage[colorlinks=true, linkcolor=cae_blue, citecolor=cae_blue, urlcolor=cae_blue]{hyperref}

% ---------- BẢNG MÀU THƯƠNG HIỆU ----------
\definecolor{cae_blue}{RGB}{0, 51, 102}
\definecolor{cae_gold}{RGB}{196, 149, 60}
\definecolor{box_bg}{RGB}{245, 247, 250}
\definecolor{box_line}{RGB}{0, 51, 102}

% ---------- TIÊU ĐỀ MỤC ----------
% Đệm đầu trang: banner letterhead chiếm 0–19mm, mà lề trên là 22mm.
% \section* đầu trang (sau \clearpage) sát lề nên phải chừa thêm, nếu không
% sẽ đè lên banner.
\titleformat{\section}
  {\normalfont\large\bfseries\sffamily\color{cae_blue}}{\thesection.}{0.6em}{}
\titleformat{\subsection}
  {\normalfont\normalsize\bfseries\sffamily\color{cae_blue!85!black}}{\thesubsection}{0.6em}{}
\titlespacing*{\section}{0pt}{14pt}{6pt}
\titlespacing*{\subsection}{0pt}{10pt}{4pt}

% Tiêu đề không đánh số ở đầu trang: \section* sau \clearpage nằm sát lề trên
% nên đè lên banner letterhead. Dùng \vspace* ngay TRƯỚC \section* không đủ,
% vì titlesec chèn khoảng cách riêng. Cách chắc chắn là đệm bằng \null.
\newcommand{\sectiondautrang}[1]{\section*{#1}}

% Dạng bài dùng \subsection*{} nên cần định dạng riêng, đậm và có gạch chân màu
\titleformat{\subsubsection}
  {\normalfont\normalsize\bfseries\sffamily\color{cae_blue!85!black}}{\thesubsubsection}{0.6em}{}

% ---------- LETTERHEAD (ảnh nền toàn trang) ----------
% Ảnh letterhead của CAE Shanghai: banner xanh trên/dưới kèm logo chìm giữa.
% Tỉ lệ ảnh 1414x2000 = 0.7070, khớp A4 (0.7071) nên chèn thẳng không méo.
% Đặt ở lớp nền dưới cùng để chữ không bị ảnh che.
% Dùng \AddToShipoutPictureBG (không có dấu *) để áp dụng cho MỌI trang.
\usepackage{eso-pic}
\AddToShipoutPictureBG{%
  \AtPageLowerLeft{%
    \includegraphics[width=\paperwidth,height=\paperheight]{images/letterhead.png}%
  }%
}

% ---------- ĐẦU / CHÂN TRANG ----------
% Letterhead đã có banner trên và dưới, nên bỏ dòng chạy để tránh chồng chữ.
% Số trang đặt trong vùng nội dung, sát trên banner dưới, lề phải.
\pagestyle{fancy}
\setlength{\headheight}{0pt}
\fancyhf{}
\renewcommand{\headrulewidth}{0pt}
\renewcommand{\footrulewidth}{0pt}
\fancyfoot[R]{\raisebox{0.9cm}{\small\color{cae_blue}\thepage}}

% ---------- MÔI TRƯỜNG NỘI DUNG ----------
% Trong tcolorbox, LaTeX đặt parindent = parskip = 0, nên các đoạn văn trong
% khung dính liền nhau. Đặt lại khoảng cách đoạn cho mọi khung.
\newcommand{\khoangcachdoan}{\setlength{\parskip}{0.5em}\setlength{\parindent}{0pt}}
\tcbset{before upper={\khoangcachdoan}}

\newtcolorbox{dinhnghia}[1][]{
  colback=box_bg, colframe=box_line, boxrule=0.6pt, arc=1pt,
  fonttitle=\bfseries, title={Định nghĩa #1}, breakable}
\newtcolorbox{dinhly}[1][]{
  colback=box_bg, colframe=cae_gold, boxrule=0.6pt, arc=1pt,
  fonttitle=\bfseries, title={Tính chất #1}, breakable}
\newtcolorbox{vidu}[1][]{
  colback=white, colframe=black!40, boxrule=0.4pt, arc=1pt,
  fonttitle=\bfseries, title={Ví dụ #1}, breakable}
\newtcolorbox{baitap}[1][]{
  colback=white, colframe=cae_blue!50, boxrule=0.4pt, arc=1pt,
  fonttitle=\bfseries, title={Câu #1}, breakable}

% Ô đáp án: viền xanh đậm hơn, dùng trong file đáp án
\newtcolorbox{dapan}[1][]{
  colback=box_bg, colframe=cae_blue, boxrule=0.6pt, arc=1pt,
  fonttitle=\bfseries, title={Câu #1}, breakable}

% ---------- TÊN CHỨNG MINH ----------
\renewcommand{\proofname}{\normalfont\bfseries Chứng minh}

% ---------- MÔN HỌC ----------
% Đổi được ở file .tex trước khi % =====================================================================
%  CAE SHANGHAI — GÓI LỆNH DÙNG CHUNG CHO TÀI LIỆU LUYỆN THI CSCA TOÁN
%  Biên dịch bằng XeLaTeX. Mọi file .tex trong buổi học \input file này.
% =====================================================================

% ---------- NGÔN NGỮ & PHÔNG CHỮ ----------
\usepackage{fontspec}
\usepackage{polyglossia}
\setmainlanguage{vietnamese}
\setotherlanguage{english}

% Chữ Latin (Anh + Việt): TeX Gyre Termes — bản clone Times New Roman,
% đi kèm TeX Live, hỗ trợ đầy đủ dấu tiếng Việt.
\setmainfont{TeX Gyre Termes}
\setsansfont{TeX Gyre Heros}

% Chữ Hán: ưu tiên Source Han Serif SC (思源宋体) — chuẩn học thuật Trung Quốc.
% Nếu máy chưa cài, tự động dùng Songti SC (宋体, có sẵn trên macOS).
% Cả hai đều thuộc họ 宋体 nên diện mạo tài liệu không đổi.
\IfFontExistsTF{Source Han Serif SC}{
  \newcommand{\cjkfamilyname}{Source Han Serif SC}
}{
  \IfFontExistsTF{Songti SC}{
    \newcommand{\cjkfamilyname}{Songti SC}
  }{
    \newcommand{\cjkfamilyname}{Noto Serif CJK SC}
  }
}

% xeCJK tự chuyển font khi gặp ký tự Hán, không cần bọc \cjk{} thủ công.
% Đây là điều kiện bắt buộc: font Latin (TeX Gyre Termes) không có glyph chữ Hán,
% nếu không khai báo xeCJK thì mọi chữ Hán viết trần sẽ bị nuốt khỏi bản in.
\usepackage{xeCJK}
\setCJKfamilyfont{caehan}{\cjkfamilyname}[Scale=0.95]
\setCJKmainfont{\cjkfamilyname}[Scale=0.95]
\setCJKsansfont{\cjkfamilyname}[Scale=0.95]
\setCJKmonofont{\cjkfamilyname}[Scale=0.95]

% \cjk{} ép dùng font Hán trong ngữ cảnh cần thiết (ví dụ trong tiêu đề in đậm
% hoặc trong bảng có định dạng riêng). Dùng \CJKfamily chứ không phải tên font,
% vì tên font chỉ là văn bản — sẽ in ra chữ thay vì đổi font.
\newcommand{\cjk}[1]{{\CJKfamily{caehan}#1}}

% ---------- TOÁN HỌC ----------
\usepackage{amsmath, amssymb, amsthm, mathtools}
\usepackage{unicode-math}
\setmathfont{TeX Gyre Termes Math}

% ---------- BỐ CỤC TRANG ----------
% Lề phải chừa vùng an toàn của letterhead:
%   banner trên chiếm 0–19mm  → top tối thiểu 2,2cm
%   banner dưới chiếm 281–297mm → bottom tối thiểu 2,2cm
% Lề trên 2,6cm (26mm) để tiêu đề đầu trang không đè banner (banner hết ở 19mm).
\usepackage[a4paper, top=2.6cm, bottom=2.4cm, left=2.5cm, right=2.0cm]{geometry}
\usepackage{fancyhdr}
\usepackage{titlesec}
\usepackage{enumitem}
\usepackage{graphicx}
\usepackage{booktabs, array, longtable, tabularx}
\usepackage{xcolor}
\usepackage[most]{tcolorbox}
\usepackage[colorlinks=true, linkcolor=cae_blue, citecolor=cae_blue, urlcolor=cae_blue]{hyperref}

% ---------- BẢNG MÀU THƯƠNG HIỆU ----------
\definecolor{cae_blue}{RGB}{0, 51, 102}
\definecolor{cae_gold}{RGB}{196, 149, 60}
\definecolor{box_bg}{RGB}{245, 247, 250}
\definecolor{box_line}{RGB}{0, 51, 102}

% ---------- TIÊU ĐỀ MỤC ----------
% Đệm đầu trang: banner letterhead chiếm 0–19mm, mà lề trên là 22mm.
% \section* đầu trang (sau \clearpage) sát lề nên phải chừa thêm, nếu không
% sẽ đè lên banner.
\titleformat{\section}
  {\normalfont\large\bfseries\sffamily\color{cae_blue}}{\thesection.}{0.6em}{}
\titleformat{\subsection}
  {\normalfont\normalsize\bfseries\sffamily\color{cae_blue!85!black}}{\thesubsection}{0.6em}{}
\titlespacing*{\section}{0pt}{14pt}{6pt}
\titlespacing*{\subsection}{0pt}{10pt}{4pt}

% Tiêu đề không đánh số ở đầu trang: \section* sau \clearpage nằm sát lề trên
% nên đè lên banner letterhead. Dùng \vspace* ngay TRƯỚC \section* không đủ,
% vì titlesec chèn khoảng cách riêng. Cách chắc chắn là đệm bằng \null.
\newcommand{\sectiondautrang}[1]{\section*{#1}}

% Dạng bài dùng \subsection*{} nên cần định dạng riêng, đậm và có gạch chân màu
\titleformat{\subsubsection}
  {\normalfont\normalsize\bfseries\sffamily\color{cae_blue!85!black}}{\thesubsubsection}{0.6em}{}

% ---------- LETTERHEAD (ảnh nền toàn trang) ----------
% Ảnh letterhead của CAE Shanghai: banner xanh trên/dưới kèm logo chìm giữa.
% Tỉ lệ ảnh 1414x2000 = 0.7070, khớp A4 (0.7071) nên chèn thẳng không méo.
% Đặt ở lớp nền dưới cùng để chữ không bị ảnh che.
% Dùng \AddToShipoutPictureBG (không có dấu *) để áp dụng cho MỌI trang.
\usepackage{eso-pic}
\AddToShipoutPictureBG{%
  \AtPageLowerLeft{%
    \includegraphics[width=\paperwidth,height=\paperheight]{images/letterhead.png}%
  }%
}

% ---------- ĐẦU / CHÂN TRANG ----------
% Letterhead đã có banner trên và dưới, nên bỏ dòng chạy để tránh chồng chữ.
% Số trang đặt trong vùng nội dung, sát trên banner dưới, lề phải.
\pagestyle{fancy}
\setlength{\headheight}{0pt}
\fancyhf{}
\renewcommand{\headrulewidth}{0pt}
\renewcommand{\footrulewidth}{0pt}
\fancyfoot[R]{\raisebox{0.9cm}{\small\color{cae_blue}\thepage}}

% ---------- MÔI TRƯỜNG NỘI DUNG ----------
% Trong tcolorbox, LaTeX đặt parindent = parskip = 0, nên các đoạn văn trong
% khung dính liền nhau. Đặt lại khoảng cách đoạn cho mọi khung.
\newcommand{\khoangcachdoan}{\setlength{\parskip}{0.5em}\setlength{\parindent}{0pt}}
\tcbset{before upper={\khoangcachdoan}}

\newtcolorbox{dinhnghia}[1][]{
  colback=box_bg, colframe=box_line, boxrule=0.6pt, arc=1pt,
  fonttitle=\bfseries, title={Định nghĩa #1}, breakable}
\newtcolorbox{dinhly}[1][]{
  colback=box_bg, colframe=cae_gold, boxrule=0.6pt, arc=1pt,
  fonttitle=\bfseries, title={Tính chất #1}, breakable}
\newtcolorbox{vidu}[1][]{
  colback=white, colframe=black!40, boxrule=0.4pt, arc=1pt,
  fonttitle=\bfseries, title={Ví dụ #1}, breakable}
\newtcolorbox{baitap}[1][]{
  colback=white, colframe=cae_blue!50, boxrule=0.4pt, arc=1pt,
  fonttitle=\bfseries, title={Câu #1}, breakable}

% Ô đáp án: viền xanh đậm hơn, dùng trong file đáp án
\newtcolorbox{dapan}[1][]{
  colback=box_bg, colframe=cae_blue, boxrule=0.6pt, arc=1pt,
  fonttitle=\bfseries, title={Câu #1}, breakable}

% ---------- TÊN CHỨNG MINH ----------
\renewcommand{\proofname}{\normalfont\bfseries Chứng minh}

% ---------- MÔN HỌC ----------
% Đổi được ở file .tex trước khi \input{cae-style} bằng \renewcommand{\monhoc}{...}
\providecommand{\monhoc}{Môn Toán}

% ---------- TRANG BÌA GỌN ----------
% #1 = loại tài liệu (TÀI LIỆU CHUẨN BỊ, TÀI LIỆU BUỔI HỌC, ...)
% #2 = số buổi
% #3 = tên buổi tiếng Việt
% #4 = tên buổi tiếng Hán
% #5 = module và tỉ trọng
\newcommand{\trangbia}[5]{%
  \begin{center}
    \vspace*{0.6cm}
    {\large\bfseries\color{cae_blue} CAE SHANGHAI}\\[0.15cm]
    {\footnotesize Tài liệu luyện thi CSCA · \monhoc}\\[0.5cm]
    {\color{cae_gold}\rule{0.55\linewidth}{0.8pt}}\\[0.35cm]
    {\small\bfseries\color{cae_blue!80!black} #1}\\[0.35cm]
    {\Large\bfseries\color{cae_blue} BUỔI #2 — #3}\\[0.2cm]
    {\large\color{cae_blue!85!black} \cjk{#4}}\\[0.25cm]
    \ifx\relax#5\relax\else{\small #5}\\[0.35cm]\fi
    {\color{cae_gold}\rule{0.55\linewidth}{0.8pt}}\\[0.3cm]
    {\footnotesize Biên soạn: Hoàng Dương, CAE Shanghai \quad·\quad \today}
  \end{center}
  \vspace{0.3cm}
}

% ---------- TRANG BÌA CHO ĐỀ KIỂM TRA ----------
% Khác \trangbia ở chỗ không có chữ "BUỔI N" — dùng cho đề thi, không phải bài giảng.
% #1 = loại tài liệu, #2 = tên kỳ thi (Việt), #3 = tên kỳ thi (Hán), #4 = thông số
\newcommand{\biade}[4]{%
  \begin{center}
    \vspace*{0.6cm}
    {\large\bfseries\color{cae_blue} CAE SHANGHAI}\\[0.15cm]
    {\footnotesize Tài liệu luyện thi CSCA · \monhoc}\\[0.5cm]
    {\color{cae_gold}\rule{0.55\linewidth}{0.8pt}}\\[0.35cm]
    {\small\bfseries\color{cae_blue!80!black} #1}\\[0.35cm]
    {\Large\bfseries\color{cae_blue} #2}\\[0.2cm]
    {\large\color{cae_blue!85!black} \cjk{#3}}\\[0.25cm]
    {\small #4}\\[0.35cm]
    {\color{cae_gold}\rule{0.55\linewidth}{0.8pt}}\\[0.3cm]
    {\footnotesize Biên soạn: Hoàng Dương, CAE Shanghai \quad·\quad \today}
  \end{center}
  \vspace{0.3cm}
}

% ---------- TIỆN ÍCH ----------
% Nhãn từ khóa in đậm ở đầu đoạn, dùng cho Từ vựng / Từ khóa nhận dạng / Chú ý
% Khoảng cách đoạn do khung tự đặt (xem \khoangcachdoan ở trên),
% ở đây chỉ cần tách đoạn và bỏ thụt đầu dòng.
\newcommand{\nhan}[1]{\par\noindent\textbf{#1.}\ \ignorespaces}

% Bảng gọn, không đánh số
\newcommand{\bang}[1]{%
  \begin{center}#1\end{center}}

% Đánh dấu ô trống cho học sinh điền
\newcommand{\chotrong}{\makebox[3.2cm]{\dotfill}}
 bằng \renewcommand{\monhoc}{...}
\providecommand{\monhoc}{Môn Toán}

% ---------- TRANG BÌA GỌN ----------
% #1 = loại tài liệu (TÀI LIỆU CHUẨN BỊ, TÀI LIỆU BUỔI HỌC, ...)
% #2 = số buổi
% #3 = tên buổi tiếng Việt
% #4 = tên buổi tiếng Hán
% #5 = module và tỉ trọng
\newcommand{\trangbia}[5]{%
  \begin{center}
    \vspace*{0.6cm}
    {\large\bfseries\color{cae_blue} CAE SHANGHAI}\\[0.15cm]
    {\footnotesize Tài liệu luyện thi CSCA · \monhoc}\\[0.5cm]
    {\color{cae_gold}\rule{0.55\linewidth}{0.8pt}}\\[0.35cm]
    {\small\bfseries\color{cae_blue!80!black} #1}\\[0.35cm]
    {\Large\bfseries\color{cae_blue} BUỔI #2 — #3}\\[0.2cm]
    {\large\color{cae_blue!85!black} \cjk{#4}}\\[0.25cm]
    \ifx\relax#5\relax\else{\small #5}\\[0.35cm]\fi
    {\color{cae_gold}\rule{0.55\linewidth}{0.8pt}}\\[0.3cm]
    {\footnotesize Biên soạn: Hoàng Dương, CAE Shanghai \quad·\quad \today}
  \end{center}
  \vspace{0.3cm}
}

% ---------- TRANG BÌA CHO ĐỀ KIỂM TRA ----------
% Khác \trangbia ở chỗ không có chữ "BUỔI N" — dùng cho đề thi, không phải bài giảng.
% #1 = loại tài liệu, #2 = tên kỳ thi (Việt), #3 = tên kỳ thi (Hán), #4 = thông số
\newcommand{\biade}[4]{%
  \begin{center}
    \vspace*{0.6cm}
    {\large\bfseries\color{cae_blue} CAE SHANGHAI}\\[0.15cm]
    {\footnotesize Tài liệu luyện thi CSCA · \monhoc}\\[0.5cm]
    {\color{cae_gold}\rule{0.55\linewidth}{0.8pt}}\\[0.35cm]
    {\small\bfseries\color{cae_blue!80!black} #1}\\[0.35cm]
    {\Large\bfseries\color{cae_blue} #2}\\[0.2cm]
    {\large\color{cae_blue!85!black} \cjk{#3}}\\[0.25cm]
    {\small #4}\\[0.35cm]
    {\color{cae_gold}\rule{0.55\linewidth}{0.8pt}}\\[0.3cm]
    {\footnotesize Biên soạn: Hoàng Dương, CAE Shanghai \quad·\quad \today}
  \end{center}
  \vspace{0.3cm}
}

% ---------- TIỆN ÍCH ----------
% Nhãn từ khóa in đậm ở đầu đoạn, dùng cho Từ vựng / Từ khóa nhận dạng / Chú ý
% Khoảng cách đoạn do khung tự đặt (xem \khoangcachdoan ở trên),
% ở đây chỉ cần tách đoạn và bỏ thụt đầu dòng.
\newcommand{\nhan}[1]{\par\noindent\textbf{#1.}\ \ignorespaces}

% Bảng gọn, không đánh số
\newcommand{\bang}[1]{%
  \begin{center}#1\end{center}}

% Đánh dấu ô trống cho học sinh điền
\newcommand{\chotrong}{\makebox[3.2cm]{\dotfill}}
 bằng \renewcommand{\monhoc}{...}
\providecommand{\monhoc}{Môn Toán}

% ---------- TRANG BÌA GỌN ----------
% #1 = loại tài liệu (TÀI LIỆU CHUẨN BỊ, TÀI LIỆU BUỔI HỌC, ...)
% #2 = số buổi
% #3 = tên buổi tiếng Việt
% #4 = tên buổi tiếng Hán
% #5 = module và tỉ trọng
\newcommand{\trangbia}[5]{%
  \begin{center}
    \vspace*{0.6cm}
    {\large\bfseries\color{cae_blue} CAE SHANGHAI}\\[0.15cm]
    {\footnotesize Tài liệu luyện thi CSCA · \monhoc}\\[0.5cm]
    {\color{cae_gold}\rule{0.55\linewidth}{0.8pt}}\\[0.35cm]
    {\small\bfseries\color{cae_blue!80!black} #1}\\[0.35cm]
    {\Large\bfseries\color{cae_blue} BUỔI #2 — #3}\\[0.2cm]
    {\large\color{cae_blue!85!black} \cjk{#4}}\\[0.25cm]
    \ifx\relax#5\relax\else{\small #5}\\[0.35cm]\fi
    {\color{cae_gold}\rule{0.55\linewidth}{0.8pt}}\\[0.3cm]
    {\footnotesize Biên soạn: Hoàng Dương, CAE Shanghai \quad·\quad \today}
  \end{center}
  \vspace{0.3cm}
}

% ---------- TRANG BÌA CHO ĐỀ KIỂM TRA ----------
% Khác \trangbia ở chỗ không có chữ "BUỔI N" — dùng cho đề thi, không phải bài giảng.
% #1 = loại tài liệu, #2 = tên kỳ thi (Việt), #3 = tên kỳ thi (Hán), #4 = thông số
\newcommand{\biade}[4]{%
  \begin{center}
    \vspace*{0.6cm}
    {\large\bfseries\color{cae_blue} CAE SHANGHAI}\\[0.15cm]
    {\footnotesize Tài liệu luyện thi CSCA · \monhoc}\\[0.5cm]
    {\color{cae_gold}\rule{0.55\linewidth}{0.8pt}}\\[0.35cm]
    {\small\bfseries\color{cae_blue!80!black} #1}\\[0.35cm]
    {\Large\bfseries\color{cae_blue} #2}\\[0.2cm]
    {\large\color{cae_blue!85!black} \cjk{#3}}\\[0.25cm]
    {\small #4}\\[0.35cm]
    {\color{cae_gold}\rule{0.55\linewidth}{0.8pt}}\\[0.3cm]
    {\footnotesize Biên soạn: Hoàng Dương, CAE Shanghai \quad·\quad \today}
  \end{center}
  \vspace{0.3cm}
}

% ---------- TIỆN ÍCH ----------
% Nhãn từ khóa in đậm ở đầu đoạn, dùng cho Từ vựng / Từ khóa nhận dạng / Chú ý
% Khoảng cách đoạn do khung tự đặt (xem \khoangcachdoan ở trên),
% ở đây chỉ cần tách đoạn và bỏ thụt đầu dòng.
\newcommand{\nhan}[1]{\par\noindent\textbf{#1.}\ \ignorespaces}

% Bảng gọn, không đánh số
\newcommand{\bang}[1]{%
  \begin{center}#1\end{center}}

% Đánh dấu ô trống cho học sinh điền
\newcommand{\chotrong}{\makebox[3.2cm]{\dotfill}}


\begin{document}

\trangbia{KẾ HOẠCH GIẢNG DẠY}{16}{Xác suất và thống kê}{概率与统计}{}

Thời lượng 90 phút (40 + 10 nghỉ + 40), thực chữa 80 phút. Tài liệu này chỉ dành cho giáo viên; đề và đáp án đã có trong \texttt{../hoc-sinh/}.

\section{PHÂN BỔ THỜI GIAN}

\bang{
\begin{tabularx}{\linewidth}{@{}Xcc@{}}
\toprule
\textbf{Hoạt động} & \textbf{Tỉ trọng} & \textbf{Thời lượng} \\
\midrule
Tổng quan và phân loại lỗi & 10\% & 8 phút \\
Chữa nhanh các câu đúng & 25\% & 20 phút \\
Chữa kỹ các câu sai & 45\% & 36 phút \\
Chiến lược xử lý nhanh nhóm câu cuối đề & 20\% & 16 phút \\
\bottomrule
\end{tabularx}}

Buổi chữa đề không giảng kiến thức mới. Mục tiêu duy nhất là biến mỗi câu sai của học sinh thành một quy tắc nhận dạng dùng lại được ở đề sau.

\section{PHÂN BỐ CÂU THEO NHÓM NỘI DUNG}

\bang{
\begin{tabularx}{\linewidth}{@{}Xlcl@{}}
\toprule
\textbf{Nhóm} & \textbf{Câu} & \textbf{Số câu} & \textbf{Ghi chú} \\
\midrule
Xác suất & 29–42 & 14 & Trọng tâm buổi, chiếm gần 30\% đề \\
Thống kê & 43–48 & 6 & Trọng tâm buổi \\
Tập hợp và bất đẳng thức & 1–7 & 7 & Ôn tích lũy \\
Hàm số và tích phân & 8–11 & 4 & Ôn tích lũy \\
Lượng giác & 12–15 & 4 & Ôn tích lũy \\
Dãy số & 16–18 & 3 & Ôn tích lũy \\
Hình học giải tích và không gian & 19–22 & 4 & Ôn tích lũy \\
Vector và số phức & 23–28 & 6 & Ôn tích lũy \\
\bottomrule
\end{tabularx}}

Đề xếp câu theo nhóm, nhóm xác suất và thống kê đặt ở cuối, đúng như thứ tự trong đề thi thật.

\section{KHỐI 1 — TỔNG QUAN VÀ PHÂN LOẠI LỖI (8 phút)}

\subsection{Mục tiêu}

Học sinh thấy được bức tranh lỗi của cả lớp trước khi chữa từng câu, để biết mình cần nghe kỹ phần nào.

\subsection{Kịch bản}

\textbf{Phút 1–3 · Trả bài và công bố phổ điểm.} Chiếu bảng phân bố đáp án, yêu cầu học sinh tự đếm số câu đúng theo từng nhóm nội dung. Học sinh ghi ra bốn con số: đúng bao nhiêu câu xác suất, bao nhiêu câu thống kê, bao nhiêu câu ôn tích lũy, và tổng thời gian đã dùng.

\textbf{Phút 4–6 · Ba nhóm lỗi của buổi này.} Trình bày ba nhóm lỗi dự kiến, không đọc từng câu:

1. \textbf{Nhầm phép cộng với phép nhân xác suất.} Đây là lỗi phổ biến nhất ở nhóm xác suất, xuất hiện ở các câu 31, 32, 39, 41.
2. \textbf{Đếm sai không gian mẫu.} Học sinh biết công thức nhưng chọn sai giữa tổ hợp và chỉnh hợp, hoặc quên rút gọn phân số về đáp án.
3. \textbf{Bỏ câu cuối vì hết giờ.} Nhóm xác suất và thống kê nằm ở cuối đề, đúng vị trí học sinh thường không còn thời gian.

\textbf{Phút 7–8 · Chốt nguyên tắc của buổi.} Nói rõ: "Nhóm xác suất và thống kê chỉ chiếm 1–2 câu trong đề thi thật, tức khoảng 3\% tổng điểm. Không được dành quá 3 phút cho cả nhóm này khi làm bài. Câu nào không ra đáp án trong 60 giây thì bỏ qua và quay lại cuối giờ."

\subsection{Điểm cần chốt}

Học sinh phải trả lời được: \textbf{"Nhóm xác suất và thống kê chiếm bao nhiêu câu trong đề thi thật?"} Đáp án là 1–2 câu. Nếu cả lớp không nhớ, nhắc lại và yêu cầu ghi vào góc vở.

\section{KHỐI 2 — CHỮA NHANH CÁC CÂU ĐÚNG (20 phút)}

\subsection{Mục tiêu}

Xác nhận cách làm đúng và rút ra quy tắc nhận dạng cho những câu học sinh làm tốt, để các em không mất thời gian suy lại ở đề sau.

\subsection{Kịch bản}

Chữa theo cụm, mỗi cụm 3–5 câu, chỉ nói \textbf{từ khóa nhận dạng} và \textbf{bước quyết định}, không trình bày lời giải đầy đủ.

\bang{
\begin{tabularx}{\linewidth}{@{}llXc@{}}
\toprule
\textbf{Cụm} & \textbf{Câu} & \textbf{Nội dung chữa} & \textbf{Thời lượng} \\
\midrule
1 & 1–7 & Tập hợp và bất đẳng thức & 5 phút \\
2 & 8–11 & Tích phân, tập xác định và nghiệm của hàm số & 4 phút \\
3 & 12–18 & Lượng giác và dãy số & 5 phút \\
4 & 19–28 & Hình học, vector và số phức & 5 phút \\
5 & 43–48 & Thống kê: trung bình, trung vị, mốt, phương sai, phân phối chuẩn & 1 phút \\
\bottomrule
\end{tabularx}}

\subsection{Từ khóa và bước quyết định theo cụm}

\textbf{Cụm 1. Tập hợp và bất đẳng thức.} Ba từ khóa quyết định: 补集 (phần bù) ở câu 1, 有且仅有一个元素 (có đúng một phần tử) ở câu 2, 不一定成立 (không nhất thiết đúng) ở câu 6. Câu 7 dùng $\sqrt{x^2-4x+4}=\lvert x-2 \rvert$ để hạ dấu căn trước khi phá dấu giá trị tuyệt đối.

\textbf{Cụm 2. Tích phân và tập xác định.} Câu 9 dùng $\displaystyle\int \dfrac{1}{x}\,dx = \ln\lvert x \rvert$. Câu 11 là bài tập xác định, phải xét đồng thời ba điều kiện: biểu thức trong căn không âm, biểu thức trong logarit dương, và mẫu khác 0.

\textbf{Cụm 3. Lượng giác và dãy số.} Câu 12 dùng $\sin^2\alpha + \cos^2\alpha = 1$ kèm điều kiện góc phần tư. Câu 14 dùng công thức nhân đôi $\sin 2\alpha = 2\sin\alpha\cos\alpha$. Câu 15 xét dấu $y' = -\sin x$ thay vì nhớ bảng biến thiên. Câu 16 và 17 chỉ cần công thức số hạng tổng quát.

\textbf{Cụm 4. Hình học, vector và số phức.} Câu 19 dùng công thức khoảng cách từ điểm đến mặt phẳng. Câu 21 nhớ quy tắc đối xứng qua gốc: đổi dấu cả ba tọa độ. Câu 22 lấy tích vô hướng hai pháp tuyến, bằng 0 thì hai mặt phẳng vuông góc. Câu 23 dùng $(1+i)^2 = 2i$ để hạ bậc lũy thừa. Câu 26 là bài nhận dạng tam giác, phải khai thác cả hai điều kiện đề cho.

\textbf{Cụm 5. Thống kê.} Ba công thức quyết định: trung bình cộng, trung vị của dãy chẵn phần tử, và phương sai. Câu 48 dùng quy tắc ba sigma: đọc $\mu$, $\sigma$ từ đề rồi so với mốc $\mu + k\sigma$.

\subsection{Câu luyện tại chỗ}

\bang{
\begin{tabularx}{\linewidth}{@{}lXl@{}}
\toprule
\textbf{Cụm} & \textbf{Câu luyện} & \textbf{Mã câu} \\
\midrule
1 & $A = \{x \in \mathbb{Z} \mid \lvert x \rvert \le 2\}$. Đếm số tập con & tự soạn \\
3 & $\cos\alpha = \dfrac{5}{13}$, $\alpha$ nhọn. Tính $\sin 2\alpha$ & tự soạn \\
4 & $z = (1 - i)^2$. Tính phần ảo & tự soạn \\
5 & Dãy $3, 5, 7, 9, 11$. Tính phương sai & tự soạn \\
\bottomrule
\end{tabularx}}

\section{KHỐI 3 — CHỮA KỸ CÁC CÂU SAI (36 phút)}

\subsection{Mục tiêu}

Mỗi câu chữa kỹ phải để lại một quy tắc viết được lên bảng và dùng lại được. Không chữa theo kiểu đọc lời giải.

\subsection{Danh sách câu cần chữa kỹ}

\bang{
\begin{tabularx}{\linewidth}{@{}cllX@{}}
\toprule
\textbf{Ưu tiên} & \textbf{Câu} & \textbf{Mã câu} & \textbf{Lý do} \\
\midrule
1 & 29 & \texttt{CAE-M-CD8-08.1-H-003} & Nhầm 合格率 (tỉ lệ đạt) với số lượng tuyệt đối \\
1 & 30 & \texttt{CAE-M-CD8-08.1-H-006} & Chia sai tỉ lệ khi làm bài toán lấy mẫu phân tầng \\
1 & 31 & \texttt{CAE-M-CD8-08.2-H-010} & Không nhận ra hai biến cố xung khắc nên cộng thẳng xác suất \\
1 & 32 & \texttt{CAE-M-CD8-08.2-H-012} & Quên trừ phần giao khi hai biến cố không xung khắc \\
1 & 33 & \texttt{CAE-M-CD8-08.3-H-018} & Giải sai phương trình $\dfrac{6}{n} = \dfrac{2}{5}$, ra số bóng đỏ thay vì số bóng xanh \\
1 & 35 & \texttt{CAE-M-CD8-08.5-H-025} & Quên dùng biến cố đối cho cụm 被选中 (được chọn) \\
1 & 36 & \texttt{CAE-M-CD8-08.5-H-026} & Không nhận ra phép lấy không hoàn lại, mẫu số phải giảm đi 1 \\
1 & 40 & \texttt{CAE-M-CD8-08.7-H-041} & Không nhận ra công thức xác suất toàn phần \\
1 & 41 & \texttt{CAE-M-CD8-08.8-H-047} & Cộng thẳng hai xác suất mà quên trừ phần giao \\
1 & 42 & \texttt{CAE-M-CD8-08.8-H-048} & Quên hệ số tổ hợp $C_{10}^9$ trong lược đồ Bernoulli \\
2 & 2 & \texttt{CAE-M-CD1-08.8-E-032} & Quên trường hợp hệ số bậc hai bằng 0 \\
2 & 11 & \texttt{CAE-M-CD4-GV-E-021} & Bỏ sót điều kiện mẫu khác 0 \\
2 & 15 & \texttt{CAE-M-CD4-GV-E-005} & Nhớ nhầm khoảng đồng biến của hàm cosin \\
2 & 26 & \texttt{CAE-M-CD7-LX2-H-038} & Chỉ dùng một trong hai điều kiện đề cho \\
3 & 43 & \texttt{CAE-M-CD9-08.4-H-013} & Nhầm mốt với giá trị lớn nhất của dãy \\
3 & 47 & \texttt{CAE-M-CD9-LX1-M-028} & Nhầm trung bình điều hòa với trung bình cộng \\
\bottomrule
\end{tabularx}}

\subsection{Kịch bản chữa câu 29 đến 42 (22 phút)}

Đây là khối nặng nhất của buổi. Chữa theo trình tự bốn bước, mỗi câu khoảng 1,5 phút:

1. \textbf{Đọc từ khóa tiếng Trung} (15 giây): yêu cầu học sinh chỉ ra từ quyết định. 概率 (xác suất), 恰好 (đúng bằng), 相互独立 (độc lập), 合格率 (tỉ lệ đạt), 分层抽样 (lấy mẫu phân tầng), 不发生 (không xảy ra), 随机抽取 (lấy ngẫu nhiên).
2. \textbf{Hỏi cách nhận dạng} (15 giây): "Đề cho biết gì và hỏi gì? Đây là bài đếm hay bài công thức?"
3. \textbf{Trình bày bước quyết định} (45 giây): chỉ viết bước then chốt, không viết toàn bộ lời giải.
4. \textbf{Chốt quy tắc} (15 giây): viết quy tắc lên góc bảng và để nguyên đến hết buổi.

\subsection{Ba quy tắc chốt ở khối này}

\bang{
\begin{tabularx}{\linewidth}{@{}llX@{}}
\toprule
\textbf{Quy tắc} & \textbf{Câu áp dụng} & \textbf{Diễn giải} \\
\midrule
Gặp 不发生 hay 不被选中 thì nghĩ ngay tới biến cố đối & 35, 37 & $P(\text{đối}) = 1 - P(\text{biến cố gốc})$, thường nhanh hơn đếm trực tiếp \\
Xung khắc thì cộng thẳng, độc lập thì trừ phần giao & 31, 32, 39, 41 & Xung khắc: $P(A \cup B) = P(A) + P(B)$. Độc lập: $P(A \cup B) = P(A) + P(B) - P(A)P(B)$ \\
Xác suất toàn phần khi đề chia hai loại sản phẩm & 40 & $P = P(\text{loại 1})P(\text{đạt} \mid \text{loại 1}) + P(\text{loại 2})P(\text{đạt} \mid \text{loại 2})$ \\
\bottomrule
\end{tabularx}}

\subsection{Kịch bản chữa câu 2, 11, 15, 26 (9 phút)}

Đây là các câu ôn tích lũy mà học sinh sai vì lý do kỹ thuật, không phải vì không biết kiến thức. Chữa nhanh, mỗi câu khoảng 2 phút, tập trung vào chỗ đặt bút sai.

\subsection{Kịch bản chữa câu 43 và 47 (5 phút)}

Hai câu này sai vì lỗi nhận dạng đại lượng, không phải lỗi tính toán. Cho học sinh tự phát hiện lỗi trước khi giáo viên chữa.

\section{KHỐI 4 — CHIẾN LƯỢC XỬ LÝ NHANH NHÓM CÂU CUỐI ĐỀ (16 phút)}

\subsection{Mục tiêu}

Học sinh có một quy trình cố định cho nhóm câu cuối đề, nhóm mà các em thường bỏ trắng vì hết giờ.

\subsection{Đặc điểm nhóm câu cuối đề}

\bang{
\begin{tabularx}{\linewidth}{@{}lX@{}}
\toprule
\textbf{Đặc điểm} & \textbf{Số liệu} \\
\midrule
Vị trí & Câu 43–48, nằm ở cuối đề \\
Số câu trong đề thi thật & 1–2 câu, khoảng 3\% tổng điểm \\
Thời gian nên dành & Tối đa 3 phút cho cả nhóm \\
Rủi ro & Đây là nhóm câu có chữ số phức tạp nhất, dễ sa vào tính toán dài \\
\bottomrule
\end{tabularx}}

\textbf{Nhấn mạnh với học sinh:} nhóm xác suất và thống kê chỉ chiếm 1–2 câu ở cuối đề thi thật. Dành quá nhiều thời gian cho nhóm này là nguyên nhân phổ biến khiến học sinh mất điểm ở những câu dễ phía trước. Câu nào không ra đáp án trong 60 giây thì bỏ qua, đánh dấu, quay lại sau khi đã làm hết phần còn lại.

\subsection{Quy trình bốn bước cho nhóm câu cuối đề}

1. \textbf{Đọc lướt đề, đánh dấu câu dễ trước.} Trong nhóm cuối đề luôn có câu chỉ cần một công thức, ví dụ tính trung bình cộng hoặc đọc tham số từ phân phối chuẩn. Làm câu đó trước.
2. \textbf{Chọn công thức trước khi tính.} Xác định đề hỏi trung bình, trung vị, mốt, phương sai hay xác suất. Mỗi đại lượng chỉ có một công thức, không cần suy luận thêm.
3. \textbf{Ước lượng đáp án rồi loại trừ.} Xác suất luôn nằm trong đoạn $[0, 1]$, nên loại ngay mọi phương án ngoài đoạn đó. Phương sai luôn không âm. Trung vị luôn nằm giữa giá trị nhỏ nhất và lớn nhất của dãy.
4. \textbf{Không giải lại từ đầu nếu không ra đáp án.} Thay đáp án vào đề để kiểm tra nhanh hơn là giải lại.

\subsection{Kỹ thuật ước lượng theo nhóm}

\bang{
\begin{tabularx}{\linewidth}{@{}lXl@{}}
\toprule
\textbf{Dạng bài} & \textbf{Kỹ thuật ước lượng} & \textbf{Câu minh họa} \\
\midrule
Xác suất cổ điển & Xác suất lớn hơn $\dfrac{1}{2}$ hay nhỏ hơn, rồi loại phương án & 31, 35, 36 \\
Biến cố đối & Kết quả phải rất gần 1 khi xác suất lẻ nhỏ & 37 \\
Biến cố độc lập & $P(A \cap B)$ phải nhỏ hơn cả $P(A)$ và $P(B)$ & 32, 39, 41 \\
Trung bình, trung vị, mốt & Đáp án phải nằm giữa giá trị nhỏ nhất và lớn nhất của dãy & 43, 45 \\
Phương sai & Phương sai là bình phương nên thường lớn hơn khoảng cách trung bình & 44 \\
Phân phối chuẩn & Đọc $\mu$ và $\sigma$ từ đề, so với mốc $\mu + k\sigma$ & 48 \\
\bottomrule
\end{tabularx}}

\subsection{Thực hành tại chỗ (6 phút)}

Phát đề rút gọn 6 câu xác suất và thống kê, yêu cầu học sinh làm trong 6 phút, sau đó tự chấm. Mục đích là để học sinh thấy với quy trình cố định thì 6 câu chỉ cần khoảng 1 phút mỗi câu.

\subsection{Điểm cần chốt}

Kết thúc buổi, học sinh phải trả lời được: \textbf{"Nếu câu xác suất không ra đáp án sau 60 giây thì làm gì?"} Đáp án là bỏ qua, đánh dấu, quay lại cuối giờ.

\section{BẢNG LỖI HAY MẮC}

\bang{
\begin{tabularx}{\linewidth}{@{}llX@{}}
\toprule
\textbf{Lỗi} & \textbf{Câu thường sai} & \textbf{Cách cảnh báo} \\
\midrule
Nhầm phép cộng với phép nhân xác suất & 31, 32, 39, 41 & Viết lên bảng: xung khắc thì \textbf{cộng thẳng}, độc lập thì \textbf{trừ phần giao}. Hai trường hợp khác nhau ở phép tính \\
Quên dùng biến cố đối cho cụm 不发生 và 不被选中 & 35, 37 & Thấy hai cụm này thì viết ngay $1 - P(\text{đối})$ trước khi làm gì khác \\
Bỏ qua trọng số khi dùng xác suất toàn phần & 40 & Đề chia sản phẩm thành hai loại thì phải nhân thêm tỉ lệ từng loại, không cộng thẳng hai tỉ lệ đạt \\
Chọn nhầm tổ hợp và chỉnh hợp & 35, 36 & Nếu thứ tự không quan trọng thì dùng $C$, có quan trọng thì dùng $P$. Đề nói 选出 (chọn ra) thì dùng $C$ \\
Quên hệ số tổ hợp trong lược đồ Bernoulli & 42 & Viết $C_n^k$ ra trước khi viết phần lũy thừa, đừng viết tắt \\
Trả lời sai đại lượng đề hỏi & 33 & Đọc kỹ chữ cuối câu hỏi: 蓝球 (bóng xanh) khác 红球 (bóng đỏ) \\
Bỏ câu cuối vì hết giờ & 43–48 & Bắt buộc làm nhóm này ở lượt hai, không để trắng \\
Chia sai mẫu số khi tính phương sai & 44 & Phương sai chia cho số phần tử, không chia cho số phần tử trừ một \\
Nhầm chu kì lượng giác & 13 & Chu kì là $\dfrac{2\pi}{\omega}$, không phải $\dfrac{\pi}{\omega}$ \\
Bỏ sót điều kiện của hàm logarit và phân thức & 11 & Viết đủ ba điều kiện trước khi kết luận tập xác định \\
Quên trường hợp hệ số bậc hai bằng 0 & 2 & Thấy tham số ở hệ số của $x^2$, xét ngay $k = 1$ \\
Chỉ dùng một điều kiện của đề & 26 & Đề cho hai đẳng thức thì phải khai thác cả hai, mỗi điều kiện cho một quan hệ \\
\bottomrule
\end{tabularx}}

\section{ĐÁP ÁN 48 CÂU}

\bang{
\begin{tabular}{@{}lcccccccccccc@{}}
\toprule
\textbf{Câu} & \textbf{1} & \textbf{2} & \textbf{3} & \textbf{4} & \textbf{5} & \textbf{6} & \textbf{7} & \textbf{8} & \textbf{9} & \textbf{10} & \textbf{11} & \textbf{12} \\
\midrule
Đáp án & B & B & A & D & A & C & C & C & A & A & A & C \\
\bottomrule
\end{tabular}}

\bang{
\begin{tabular}{@{}lcccccccccccc@{}}
\toprule
\textbf{Câu} & \textbf{13} & \textbf{14} & \textbf{15} & \textbf{16} & \textbf{17} & \textbf{18} & \textbf{19} & \textbf{20} & \textbf{21} & \textbf{22} & \textbf{23} & \textbf{24} \\
\midrule
Đáp án & D & A & C & B & D & D & D & D & C & B & D & B \\
\bottomrule
\end{tabular}}

\bang{
\begin{tabular}{@{}lcccccccccccc@{}}
\toprule
\textbf{Câu} & \textbf{25} & \textbf{26} & \textbf{27} & \textbf{28} & \textbf{29} & \textbf{30} & \textbf{31} & \textbf{32} & \textbf{33} & \textbf{34} & \textbf{35} & \textbf{36} \\
\midrule
Đáp án & A & D & C & C & A & A & C & D & D & B & A & A \\
\bottomrule
\end{tabular}}

\bang{
\begin{tabular}{@{}lcccccccccccc@{}}
\toprule
\textbf{Câu} & \textbf{37} & \textbf{38} & \textbf{39} & \textbf{40} & \textbf{41} & \textbf{42} & \textbf{43} & \textbf{44} & \textbf{45} & \textbf{46} & \textbf{47} & \textbf{48} \\
\midrule
Đáp án & D & C & B & B & B & B & A & C & B & B & C & D \\
\bottomrule
\end{tabular}}

Phân bố: A 12 câu, B 12 câu, C 12 câu, D 12 câu.

\textbf{Câu cần chữa kỹ nếu học sinh làm sai nhiều:} Câu 31, 32, 35, 40, 41, 42.

\begin{itemize}[leftmargin=1.4em, itemsep=2pt]
\item \textbf{Câu 31} (\texttt{CAE-M-CD8-08.2-H-010}): học sinh dùng phép nhân cho hai biến cố xung khắc, chọn 0.02 thay vì 0.03.
\item \textbf{Câu 32} (\texttt{CAE-M-CD8-08.2-H-012}): cộng thẳng $0.4 + 0.5$ mà quên trừ $P(A \cap B) = 0.2$, chọn 0.9.
\item \textbf{Câu 35} (\texttt{CAE-M-CD8-08.5-H-025}): học sinh đếm $C_5^2 = 10$ nhưng không nghĩ tới biến cố đối, dễ đếm nhầm số cặp không chứa A.
\item \textbf{Câu 40} (\texttt{CAE-M-CD8-08.7-H-041}): cộng hai tỉ lệ đạt mà quên nhân với tỉ lệ từng loại sản phẩm.
\item \textbf{Câu 41} (\texttt{CAE-M-CD8-08.8-H-047}): cộng thẳng hai xác suất mà quên trừ phần giao, chọn 0.8.
\item \textbf{Câu 42} (\texttt{CAE-M-CD8-08.8-H-048}): chọn phương án thiếu hệ số tổ hợp $C_{10}^9$.
\end{itemize}

\section{GHI CHÚ SAU BUỔI HỌC}

\emph{(Điền sau khi dạy xong, dùng để điều chỉnh buổi sau và các buổi lặp lại.)}

\bang{
\begin{tabularx}{\linewidth}{@{}Xl@{}}
\toprule
\textbf{Nội dung} & \textbf{Ghi nhận} \\
\midrule
Thời gian thực tế từng khối &  \\
Câu học sinh sai nhiều nhất &  \\
Câu hỏi học sinh hỏi nhiều &  \\
Điều chỉnh cho lần dạy sau &  \\
\bottomrule
\end{tabularx}}

\emph{Tài liệu nội bộ, CAE SHANGHAI.}
\vspace{0.8cm}
\begin{center}
{\footnotesize\color{cae_blue} Tài liệu nội bộ, CAE SHANGHAI}
\end{center}

\end{document}
